\documentclass{article}

% Language, fonts and typography
\usepackage[british]{babel}
\usepackage[T1]{fontenc}
\usepackage{lmodern}
\usepackage{microtype}

% Mathematics
\usepackage{amsmath}
\usepackage{amssymb}
\usepackage{amsthm}

% Layout
\usepackage{fullpage}
\usepackage{parskip}
\usepackage{fancyhdr}
\usepackage{enumitem}

% Cross-references: hyperref and cleveref are loaded last, in this order
\usepackage{xcolor}
\usepackage{hyperref}
\usepackage[capitalise,nameinlink]{cleveref}

\hypersetup{
  colorlinks=true,
  allcolors=blue!50!black,
  pdftitle={The Galois extension argument},
  pdfauthor={Joe Lamond and Joshua Munchau},
}

% Page style
\pagestyle{fancy}
\fancyhf{}
\fancyhead[L]{The Galois extension argument}
\fancyhead[R]{Joe Lamond and Joshua Munchau}
\fancyfoot[C]{\thepage}
\setlength{\headheight}{14pt}
\setlength{\headsep}{10pt}

% Theorems. With parskip, the default amsthm style spaces a theorem like an ordinary
% paragraph, so it does not stand out. These values (added to parskip's paragraph gap) give
% about 14pt of white space on each side, matching the space before a proof.
\newtheoremstyle{spacedplain}{6pt}{3pt}{\itshape}{}{\bfseries}{.}{.5em}{}
\newtheoremstyle{spacedremark}{6pt}{3pt}{\normalfont}{}{\bfseries}{.}{.5em}{}
\theoremstyle{spacedplain}
\newtheorem{theorem}{Theorem}
\theoremstyle{spacedremark}
\newtheorem*{remark}{Remark}

\DeclareMathOperator{\Gal}{Gal}
\DeclareMathOperator{\Aut}{Aut}

\begin{document}

\section*{The Galois extension argument}

This is Step~4 of the proof sketch in the \emph{Outline of formalisation}.

\begin{theorem}\label{thm:galois-extension}
Suppose $R$ is a real closed field, and $C$ is an extension of $R$ for which
there is an element $i\in C$ such that $i^2=-1$ and $C=R(i)$. If $R\subseteq
C\subseteq L$ is a tower of field extensions such that $R\subseteq L$ is finite
Galois, then $L=C$.
\end{theorem}

The proof relies on two of the earlier steps of the proof sketch:
\begin{enumerate}[label=\textbf{Step \arabic*.}, start=2, leftmargin=*]
\item $C$ has no quadratic extension.
\item Every finite Galois extension of $R$ has degree a power of $2$.
\end{enumerate}

\begin{proof}
Since $R\subseteq C\subseteq L$ and $R\subseteq L$ is finite Galois, by standard
results from Galois theory it follows that $C\subseteq L$ is finite Galois, and
hence the Galois group $H=\Aut_{C}(L)$ has order $[C:L]$. By Step~3, $[L:R]=2^n$
for some $n\in\mathbb N$. The tower law then tells us that $[L:R]=[L:C][C:R]$,
so that $[L:C]=2^m$ for some $m\in\mathbb N$. If $m=0$ then we are done, so let
us assume for contradiction that $m\ge 1$. 

Recall that if $p$ is prime, $s\in\mathbb N$, and $G$ is a finite group of order $p^s$, then $G$
has a subgroup of order $p^t$ for each integer $t\in\{0,\dots,s\}$. In
particular, $H$ has a subgroup $K\subseteq H$ of index $2$. Put $F=L^K$. By
the Galois correspondence for $C\subseteq L$, we have $K=\Aut_F(L)$. Since
$F\subseteq L$ is finite Galois,
$$
[L:F]=|\Aut_F(L)|=|K|=2^{m-1} \, .
$$
By the tower law again, $[L:C]=[L:F][F:C]$, hence $[F:C]=2$. This contradicts
the fact that $C$ has no quadratic extension (Step~2).
\end{proof}

\end{document}
